% page 237 du fac-simile (image p-236 du scan)
% bloc 165.1 x 229.0 mm ; marge gauche 36.9 mm
\pgc{15.240mm}{0.000mm}{
\l{\cel{56}229}
\l{}
\l{}
\l{\sou{imaj}o. I. Aparajo videbla ek objekto lumoza, konsistanta ye la}
\l{\cel{2}radii lumala qui spricas de olca. - II. Aparajo videbla ek}
\l{\cel{2}objekto, imitita per desegnuro, pikturo, skulturo. - III.}
\l{\cel{2}Aparajo videbla di objekto, konceptita per la imagino-fakul-}
\l{\cel{2}tato. - EFIS.}
\l{}
\l{\sou{imana}r. (\sou{netrans.}) Efikar konstante, vice efikar provizore}
\l{\cel{2}(\sou{aludante Deo}).\cel{2}DEFIS.}
\l{}
\l{\sou{imbasta}r. (\sou{trans.})\cel{2}Sutar per longa steburi ed (ordinare) per}
\l{\cel{2}filo blanka, por pose probar (la vesto, la robo, e c.) - EFI.}
\l{}
\l{\sou{imbecil}a. Di qua la spirito esas febla. - DEFIS.}
\l{}
\l{huibar. (\sou{trans., ulo, ye)}. Humidigar, per igar absorbar liquido}
\l{\cel{2}maxima-quante. - DEFIS.}
\l{}
\l{\sou{imbrika}r.(\sou{trans.})\cel{2}Dispozar, aranjar quale la teguli di tekto.}
\l{\cel{2}EFIS.}
\l{}
\l{\sou{imens}a.Di qua la grandeso-grado impedas la mezuro. - EF.}
\l{}
\l{\sou{imersa}r. (\sou{trans., aden}). Parte sinkar (korpo) aden liquido.}
\l{\cel{2}- DEFIS.}
\l{}
\l{\sou{imita}r.(\sou{trans.}) I. Agar quale ulu ; facar (ulo) segun, (kom}
\l{\cel{2}modelo), la verko da ulu. - DEFIRS.}
\l{}
\l{\sou{imobl}o.(\sou{yuro-cienco}).\cel{2}Havajo nemovebla (domo, agro, e c.).DEFIS.}
\l{}
\l{\sou{imola}r.(\sou{trans.})\cel{2}I. Ocidar sakrifike por satisfacar la deajo.}
\l{\cel{2}II. Perisigar, ruinar, ulu por satisfacar interesto, pasiono,}
\l{\cel{2}e c. - EFI.}
\l{}
\l{\sou{imortel}o. (\sou{bot.}) Planto ek la familio \textquotedbl{}kompozaji\textquotedbl{}, di qua la}
\l{\cel{2}floro ne velkas, ed uzata por ornar la tombi. - L.\cel{1}\sou{helichrysum}.}
\l{\cel{2}DEFR.}
\l{}
\l{\sou{impeda}r. (\sou{trans., pri, per}).Entravar, embarasar (ulu) pri lua}
\l{\cel{2}ago od agado ; entravar la evento (di ulo). - EFIS.}
\l{}
\l{\sou{impera}r. (\sou{trans.}) Dicar (ad ulu) ke lu agez (to o co), kun im-}
\l{\cel{2}plicita nuanco di puniso kaze di neobedio. (Ne uzesas pri la}
\l{\cel{2}armeo-autoritato. Komp. \textquotedbl{}komandar\textquotedbl{}). - EFIS.}
\l{}
\l{\sou{imperativ}o. I. (\sou{gram}.)Modo qua expresas la impero. - II. (\sou{filoz.})}
\l{\cel{2}Ago-regulo absoluta e ne-kondicionoza. - DEFIS.}
\l{}
\l{imperatoro.(\sou{historio})\cel{2}Komandanto di la armeo Romana antiqua.}
\l{\cel{2}- DEFIRS.}
}
